\documentclass[]{article}

\usepackage{amsmath}
\usepackage{multirow}
\usepackage{natbib}
\usepackage{graphicx} 


\begin{document}

Turbulent flows, ubiquitous from industrial processes to astrophysical phenomena, are fundamentally governed by the dynamics of coherent structures. Among these, vortices—regions of intense, swirling motion—are not merely passive features but act as the primary agents for the transport of energy, momentum, and scalar quantities across a vast range of scales. The intermittent and non-Gaussian statistics of turbulence are intrinsically linked to the life cycle of these filament-like structures. A quantitative understanding of their Lagrangian motion, tracking how they are advected and distorted by the surrounding flow, is therefore essential for advancing the fundamental theory of turbulence and for developing more physically grounded sub-grid scale models for large-scale simulations.A central open question concerns the statistical nature of this transport. The motion of a vortex arises from a complex interplay between its own coherent, self-sustaining dynamics and its advection by the chaotic, multiscale velocity field in which it is embedded. This raises a fundamental question: does the vortex behave like a passive tracer, executing a random walk akin to classical Brownian motion, or does its structural integrity and inertia induce temporal correlations in its trajectory  \citep{maiti2025vorticitypackingeffectslong}? The answer has profound implications for modeling turbulent mixing and dissipation. A simple random walk leads to classical diffusion, where the Mean Squared Displacement (MSD) of particles grows linearly with time, $\langle |\mathbf{r}(t+\tau) - \mathbf{r}(t)|^2 \rangle \propto \tau$. However, if the motion is correlated, it can result in anomalous diffusion, such as superdiffusion where the MSD scales as $\langle |\mathbf{r}(t+\tau) - \mathbf{r}(t)|^2 \rangle \propto \tau^\alpha$ with $\alpha > 1$, indicating a far more efficient transport mechanism than classical models predict  \citep{leoncini2002jetsstickinessanomaloustransport,maiti2025vorticitypackingeffectslong}.In this work, we investigate this question by performing a detailed statistical analysis of vortex trajectories extracted from a high-resolution direct numerical simulation of three-dimensional, homogeneous, and isothermal turbulence. We utilize a robust computational pipeline to first identify coherent vortex structures using the Q-criterion, which isolates regions where rotation dominates strain. Subsequently, we track the vorticity-weighted centroids of these structures through time, generating a comprehensive dataset of thousands of Lagrangian trajectories. This direct, Lagrangian perspective allows us to probe the dynamics of the fundamental building blocks of the turbulent flow, moving beyond the limitations of static, Eulerian field descriptions.To build a self-consistent physical picture of vortex transport, we employ a multi-faceted statistical framework  \citep{kumar2025trappingtransportinertialparticles,aulnette2025transportsphericalmicroparticles3d}. First, we compute the Mean Squared Displacement to directly measure the diffusive exponent and establish the overall transport regime. To uncover the physical mechanism behind the observed behavior, we then calculate the Velocity Autocorrelation Function, which quantifies the persistence or 'memory' of a vortex's velocity. We further probe the underlying stochastic process by analyzing the probability distribution of trajectory step sizes, testing for deviations from the Gaussian statistics of a simple random walk towards heavy-tailed distributions characteristic of Lévy flights  \citep{kumar2025trappingtransportinertialparticles}. Finally, we connect these statistical properties to the vortex's physical orientation by examining the anisotropy of its motion relative to the local vorticity vector. By synthesizing these complementary analyses, we aim to determine whether Lagrangian vortex transport is best described as a correlated random walk and to quantify the degree to which it deviates from classical diffusion  \citep{kumar2025trappingtransportinertialparticles,aulnette2025transportsphericalmicroparticles3d}.

\end{document}
                